\documentclass[aspectratio=169,notheorems]{corvinusmetropolis}
\newtheorem{theorem}{{Tétel.}}
\newtheorem{proposition}[theorem]{{Állítás.}}
\newtheorem{lemma}[theorem]{{Lemma.}}
\newtheorem{FA}[theorem]{{Az algebra alaptétele.}}
\newtheorem{corollary}[theorem]{{Következmény.}}
\theoremstyle{definition}
\newtheorem{definition}[theorem]{{Definíció.}}
\newtheorem{defprop}[theorem]{{Definíció-Állítás.}}
\newtheorem{notes}[theorem]{{Megjegyzés.}}

\email{magyarkuti@uni-corvinus.hu}
\author[Magyarkuti]{Gyula~Magyarkuti\,\orcidlink{0000-0003-4444-0481}}
\institute[Corvinus Egyetem]%
{%\inst{1}
  Budapesti Corvinus Egyetem\\
  Matematika Tanszék\\
  Sóház, 208/B

  \begin{minipage}[b]{0.8\linewidth}
  email: \putemail\\
  letölthető: \url{https://magyarkuti.github.io/algebra-1}
  \end{minipage}
  \hfill{\adjustbox{valign=B}{\qrcode[hyperlink]{https://magyarkuti.github.io/algebra-1}}}%% iff qrcode needed, otherwise comment this line!
}
\date[Tavasz]{\ontoday}%% overwrite later, if needed.

\titlegraphic{\hfill{\includegraphics[width=.2\linewidth]{\myinstlogo}}}

\renewcommand{\emph}{\alert}
\renewcommand{\mathbf}{\mathbb}
\newcommand{\lin}{\operatorname{lin}}
\newcommand{\gen}{\operatorname{gen}}
\newcommand{\im}{\operatorname{Im}}
\newcommand{\id}{\operatorname{id}}
\newcommand{\sgn}{\operatorname{sgn}}
\newcommand{\tr}{\operatorname{tr}}
\newcommand{\codim}{\operatorname{codim}}
%newcommand{\deg}{\operatorname{deg}}
\newcommand{\ip}[2]{\langle#1,#2\rangle}
\newcommand{\uj}{\text{új}}
\newcommand{\rgi}{\text{régi}}
\newcommand{\Star}[1]{#1\ensuremath{^*}\kern-\scriptspace}
\newcommand{\CStar}{\Star{\ensuremath{\mathrm{C}}}}
\endinput

\title[Algebra 1.]{Mátrixok}
\subtitle{Komplex számok mátrix reprezentációja}
\date[5. hét]{2026.~október 12-én}
\begin{document}
\openingframes{Tartalom}
\section{Mátrixok}
\frame{
\begin{definition}[mátrix]
	Egy tetszőleges test feletti mátrixnak nevezzünk,
	a test elemeiből képzett táblázatot.
	Ha $m,n\in\mathbb{N}$ előre rögzített pozitív egészek és
	az $A$ táblázatnak $m$ sora és $n$ oszlopa van,
	akkor azt mondjuk, hogy $A$ egy $m\times n$ méretű mátrix.
	Az $\mathbb{F}$ test feletti $m\times n$-es mátrixok halmazát $\mathbb{F}^{m\times n}$
	módon jelöljük.

	Ha $A\in\mathbb{F}^{m\times n}$ egy mátrix,
	akkor $A_i$ jelöli az $i$-edik sort, ami persze egy $1\times n$-es mátrix;
	$A^j$ jelöli a $j$-edik oszlopot, ami persze egy $m\times 1$-s mátrix;
	$A_i^j$ jelöli az $i$-edik sor $j$-edik elemét.
	Sokszor használjuk az $A_{i,j}=A_i^j$ jelölést is.

	Időnként, azt hangsúlyozandó hogy mátrixokról van szó a mátrixot jelölő betűt kapcsos zárójelbe teszem.
	Pl. $\left[ A \right]\in\mathbb{F}^{m\times n}$.

	A mátrixot a mérete és az elemei határozzák meg.
	Emiatt két mátrix akkor azonos, ha azonos méretűek, és a megfelelő elemeik is azonosak.
\end{definition}
}
\frame{
Az azonos típusú mátrixok közt műveleteket definiálunk:
\begin{definition}[mátrixok összege]
	Rögzített $m,n\in\mathbb{N}$ mellett, ha $A,B\in\mathbb{F}^{m\times n}$,
	akkor ezek összege az a $C\in\mathbb{F}^{m\times n}$ mátrix, amelyre
	\[
		C_{i,j}=A_{i,j}+B_{i,j}
	\]
	minden $i=1,\dots,m$ és $j=1,\dots, n$.
	Jelölés: $C=A+B$.
    \index{matrix@mátrixok összege}
\end{definition}
\begin{proposition}
	Az $m\times n$ méretű mátrixok az fent definiált összeadás művelettel
	Abel-csoportot alkotnak.
	A $[0]$-val jelölt neutrális elem az az $m\times n$-s mátrix,
	amelynek minden eleme a test zérus eleme:
	\[
		[0]_{i,j}=0;
	\]
	az $[A]$ mátrix additív inverze az az $[-A]$-val jelölt $m\times n$ méretű mátrix,
	amelyre
	\[
		[-A]_{i,j}=-([A]_{i,j}).\qedhere
	\]
    \label{pr:matrixokVS1}
\end{proposition}
}
\frame{
Most definiáljuk egy számnak és egy mátrixnak a szorzatát.
\begin{definition}[szám és mátrix szorzata]
	Ha $\alpha\in\mathbb{F}$ egy szám és $A\in\mathbb{F}^{m\times n}$ egy mátrix, akkor ezek szorzata
	az $\alpha A\in\mathbb{F}^{m\times n}$ módon jelölt mátrix, melynek elemeire
	\[
		[\alpha A]_{i,j}=\alpha[A]_{i,j}.\qedhere
	\]
    \index{szám és mátrix szorzata}
\end{definition}
Könnyen ellenőrizhetőek a következő számolási szabályok:
\begin{proposition}
	Legyenek $A,B\in\mathbb{F}^{m\times n}$ mátrixok, és $\alpha,\beta\in\mathbb{F}$ tetszőleges számok.
	Ekkor
	\begin{enumerate}
		\item $\alpha\left( A+B \right)=\alpha A+\alpha B$;
		\item $\left( \alpha+\beta \right)A=\alpha A+\beta A$;
		\item $\left( \alpha\beta \right)A=\alpha\left( \beta A \right)$;
		\item $1 A=A$.\qedhere
	\end{enumerate}
    \label{pr:matrixokVS2}
\end{proposition}
}
\frame{
Az utolsó két állítást, \ref{pr:matrixokVS1}. és \ref{pr:matrixokVS2}.,
együtt később úgy fogjuk fogalmazni,
hogy adott test feletti tetszőleges méretű mátrixok \emph{vektorteret}\index{vektortér} alkotnak.
Ha az $A_1,\dots,A_k$ azonos méretű mátrixok adottak,
akkor tetszőleges $\alpha_1,\dots,\alpha_k\in\mathbb{F}$
számok mellett értelmes az $A_1,\dots,A_k$ mátrixoknak az $\alpha_1,\dots,\alpha_k$ számokkal
mint együtthatókkal képzett \emph{lineáris kombinációja}\index{lineáris kombináció}:
\[
	\alpha_1A_1+\alpha_2A_2+\ldots+\alpha_kA_k.
\]
Minden eddig megértett dolog persze az $m=1$ speciális esetben,
és az $n=1$ speciális esetben is igaz.
Ha $m=1$, akkor a mátrixot \emph{sorvektornak}\index{sorvektor};
ha pedig $n=1$, akkor a mátrixot \emph{oszlopvektornak}\index{oszlopvektor} mondjuk.
}
\frame{
Most a mátrixok között a szorzás műveletet definiáljuk.
\begin{definition}[mátrixok szorzata]\index{matrix@mátrixok szorzata}
	Legyen $A\in\mathbb{F}^{m\times k}$ és $B\in\mathbb{F}^{k\times n}$ mátrix.
	Fontos, hogy $A$ oszlopainak száma azonos $B$ sorainak számával.
	Ezek $C=AB$ szorzata egy $C\in\mathbb{F}^{m\times n}$ mátrix,
	melynek elemeit az alábbi egyenlőség definiálja
	\[
		[C]_{i,j}=\sum_{s=1}^k[A]_{i,s}[B]_{s,j}.
	\]
	Itt persze $i=1,\dots,m$ és $j=1,\dots,n.$
\end{definition}
}
\frame{
Figyeljünk arra, hogy a fenti definícióban az $A$ és $B$ mátrixok sorrendje is fontos.
Például, ha $n\neq m$, akkor $BA$ nem is értelmes.
Még $m=n$ esetén is $BA$ egy $k\times k$ méretű mátrix, míg $AB$ egy $n\times n$ méretű mátrix,
azaz az $AB$ szorzatnak általában véve semmi köze a $BA$ szorzathoz.

A bevezetett összeadás és szorzás műveletekre teljesül a disztributivitás is:
\begin{proposition}
	Legyen $A\in\mathbb{F}^{m\times k}$, valamint a $B,C\in\mathbb{F}^{k\times n}$ mátrixok.
	Ekkor
	\begin{displaymath}
		A\left( B+C \right)=AB+AC.
	\end{displaymath}
	Hasonlóan,
	ha $A,B\in\mathbb{F}^{m\times k}$, valamint a $C\in\mathbb{F}^{k\times n}$ mátrixok,
	akkor
	\begin{displaymath}
		\left( A+B \right)C=AC+BC.\qedhere
	\end{displaymath}
\end{proposition}

%\begin{proof}
%	Az első disztributivitást mutatjuk meg, a második evvel analóg.
%	Világos, hogy mindkét oldalon azonos méretű mátrixok vannak.
%	Minden szóba jövő $i,j$ index mellett
%	\begin{multline*}
%		\left[ A\left( B+C \right) \right]_{i,j}
%		=
%		\sum_{s=1}^k[A]_{i,s}\left[ B+C \right]_{s,j}
%		=
%		\sum_{s=1}^k[A]_{i,s}\left([B]_{s,j}+[C]_{s,j} \right)
%		=
%		\sum_{s=1}^k\left( [A]_{i,s}[B]_{s,j}\right)+\left([A]_{i,s}[C]_{s,j} \right)
%		=
%		\\
%		\sum_{s=1}^k[A]_{i,s}[B]_{s,j}
%		+
%		\sum_{s=1}^k[A]_{i,s}[C]_{s,j}
%		=
%		\left[ AB \right]_{i,j}+\left[ AC \right]_{i,j}
%		=
%		[AB+AC]_{i,j}.
%	\end{multline*}
%	Ez éppen  az $A\left( B+C \right)=AB+AC$ mátrixazonosságot jelenti.
%\end{proof}
}
\frame{
Szinte banalitás,
de egy mátrix az elemei összesége, de a sorai összesége, sőt az oszlopai összesége is.

A következő ultrafontos állítás arról szól, hogy a szorzás operációt hogyan kell látnunk attól függően,
hogy a szorzat mátrixra,
mint elemei összeségére,
mint oszlopai összeségére,
vagy mint sorainak összeségére gondolunk.
}
\frame{
	Legyen $A\in\mathbb{F}^{m\times k}$ és $B\in\mathbb{F}^{k\times n}$ mátrix.
	Jelölje $C=AB$ ezek szorzatát ebben a sorrendben.
	Ekkor
	\begin{itemize}
		\item A szorzat mátrix $i$-edik sorának $j$-edik eleme,
		      az $A$ mátrix $i$-edik sorának és a $B$ mátrix $j$-edik
		      oszlopának, mint speciális mátrixoknak a szorzata.
              \\
		      Magyarul: minden $1\leq i\leq m$ és $1\leq j \leq n$ mellett
		      \[
			      [C]_{i,j}=[A]_i\cdot [B]^j.
		      \]
		\item
		      A szorzat mátrix minden oszlopa az $A$ mátrix oszlopainak a $B$ mátrix megfelelő oszlopából vett elemekkel képzett
		      lineáris kombinációja.
              \\
		      Magyarul: minden $1\leq j\leq n$ mellett
		      \[
			      [C]^j=\sum_{s=1}^k[B]_s^j[A]^s.
		      \]
    \end{itemize}
}\frame{
\begin{itemize}
		\item
		      A szorzat mátrix minden sora a $B$ mátrix sorainak az $A$ mátrix megfelelő sorából vett elemekkel képzett
		      lineáris kombinációja.
              \\
		      Magyarul: minden $1\leq i\leq m$ mellett
		      \[
			      [C]_i=\sum_{s=1}^k[A]_i^s[B]_s.
		      \]
		\item
		      A szorzat mátrix az $A$ oszlopaiból, és a $B$ soraiból alkotott diádok összege.\index{diad@diád}
              \\
		      Magyarul:
		      \[
			      [C]=\sum_{s=1}^k[A]^s[B]_s.
		      \]
		      \emph{Diádnak}\index{diad@diád} nevezzük egy oszlop- és egy sorvektor szorzatát.
		      Ha az oszlopnak és a sornak rendre azonosak az elemei, akkor
		      \emph{szimmetrikus diádról} beszélünk.\index{szimmetrikus diád}
		      \qedhere
	\end{itemize}
%\begin{proof}[1. bizonyítása]
%	$[A]_i\cdot [B]^j=\sum_{s=1}^k[A]_{i,s}[B]_{s,j}=[C]_{i,j}.$
%\end{proof}
%\begin{proof}[2. bizonyítása]
%	\(
%	\left[ \sum_{s=1}^k[B]_s^j[A]^s \right]_i=
%	\sum_{s=1}^k[B]_s^j[A]_i^s =
%	\sum_{s=1}^k[A]_i^s[B]_s^j =
%	[C]_{i,j}=
%		[C]_i^j
%	\)
%	minden $i$-re.
%\end{proof}
%\begin{proof}[3. bizonyítása]
%	\(
%	\left[ \sum_{s=1}^k[A]_i^s[B]_s \right]^j=
%	\sum_{s=1}^k[A]_i^s[B]_s^j=
%	[C]_{i,j}=
%		[C]_i^j
%	\)
%	minden $j$-re.
%\end{proof}
%\begin{proof}[4. bizonyítása]
%	\(
%	\left[ \sum_{s=1}^k[A]^s[B]_s \right]_{i,j}=
%	\sum_{s=1}^k\left[ [A]^s[B]_s \right]_{i,j}=
%	\sum_{s=1}^k [A]_i^s[B]_s^j=
%	C_{i,j}
%	\)
%	minden $i$-re $j$-re.
%\end{proof}
}
\frame{
\begin{definition}[Kronecker-delta, identitás mátrix]\index{Kronceker-delta}\index{identitás mátrix}
	\emph{Kroencker-deltának} nevezzük az alábbi egyszerű szimbólumot:
	\[
		\delta_{i,j}=
		\begin{cases}
			1,\text{ ha }i=j; \\
			0,\text{ egyébként.}
		\end{cases}
	\]
	Adott $n\geq 1$ természetes számra az $n\times n$ méretű identitás mátrix azaz $I\in\mathbb{F}^{n\times n}$ mátrix,
	amelyre
	\[
		[I]_{i,j}=\delta_{i,j}.\qedhere
	\]
\end{definition}
Nyilvánvaló, hogy ha $A\in\mathbb{F}^{m\times n}$ mátrix és $I\in\mathbb{F}^{n\times n}$ méretű identitás mátrix,
akkor $A\cdot I=A$. Hasonlóan, ha most $I$ az $m\times m$ identitás mátrixot jelöli, akkor pedig $I\cdot A=A$ azonosság teljesül.
Persze, ha $A\in\mathbb{F}^{n\times n}$ négyzetes mátrix és $I\in\mathbb{F}^{n\times n}$ az ugyanilyen méretű identitás mátrix,
akkor
\[
	I\cdot A=A\cdot I=A
\]
is fennáll.
}
\frame{
A mátrixok szorzásának legérdekesebb tulajdonsága a szorzás asszociativitása.
\begin{proposition}
	Legyen az $A,B,C$ mátrixok úgy megadva, hogy $AB$ is értelmes és $BC$ is értelmes legyen,
	azaz $A\in\mathbb{F}^{m\times k},B\in\mathbb{F}^{k\times l}, C\in\mathbb{F}^{l\times n}$.
	Ekkor
	\[
		A\left( BC \right)=\left( AB \right)C.\qedhere
	\]
\end{proposition}
%\begin{proof}
%	Először is azt vegyük észre,
%	hogy ha az $A$ és a $C$ mátrixok egyike egy szám,
%	és a másik két mátrix összeszorozható,
%	akkor a mátrix szorzás definíciója szerint az állítás nyilvánvaló.
%
%	Másodszor azt vegyük észre, hogy mindkét oldalon azonos méretű
%	konkrétan $m\times n$ méretű mátrixok szerepelnek.
%
%	Azt kell tehát még meggondolnunk,
%	hogy az $i$-edik sor $j$-edik eleme mindkét oldalon ugyanaz.
%	A jobboldalon ez
%	\begin{multline*}
%		[AB]_i\cdot [C]^j
%		=
%		\left( \sum_{s=1}^k[A]_i^s[B]_s \right)[C]^j
%		=
%		\sum_{s=1}^k\left([A]_i^s[B]_s\right)[C]^j
%		=
%		\sum_{s=1}^k[A]_i^s\left([B]_s[C]^j\right)
%		\\
%		=
%		\sum_{s=1}^k[A]_i^s\left(\sum_{r=1}^l[B]_s^r[C]_r^j\right)
%		=
%		\sum_{s=1}^k\sum_{r=1}^l[A]_i^s([B]_s^r[C]_r^j).
%	\end{multline*}
%	A baloldalon
%	az $i$-edik sor $j$-edik eleme hasonló számolgatással:
%	\begin{multline*}
%		[A]_i[BC]^j=
%		[A]_i\left( \sum_{r=1}^l[B]^r[C]_r^j \right)
%		=
%		\sum_{r=1}^l[A]_i\left([B]^r[C]_r^j\right)
%		=
%		\sum_{r=1}^l\left([A]_i[B]^r\right)[C]_r^j
%		\\
%		\sum_{r=1}^l\left(\sum_{s=1}^k[A]_i^s[B]_s^r\right)[C]_r^j
%		=
%		\sum_{r=1}^l\sum_{s=1}^k\left([A]_i^s[B]_s^r\right)[C]_r^j.
%	\end{multline*}
%	A testben fennálló asszociativitás és kommutativitás miatt a bal- és a jobboldali kifejezés azonos.
%\end{proof}
}
\frame{
Ha tehát adottak az $A_1,A_2,\dots,A_p$ mátrixok, úgy hogy a sorban egymás után következőknek értelmes a szorzata,
akkor ezek bármilyen zárójelezésével képzett szorzata is értelmes,
és ugyanazt a mátrixot eredményezi.
Mint azt a számok esetén megszoktuk,
használhatjuk az $A_1\cdot A_2\cdot\ldots\cdot A_p$ jelölést ezen mátrixok bármelyik zárójelezésével képzett szorzatára.
Érdemes a fenti bizonyításból megjegyezni a formulát,
amelyet három mátrix szorzatára kapunk:
\[
	\left[ ABC \right]_{i,j}
	=
	\sum_{s=1}^k\sum_{r=1}^l\left[ A \right]_{i,s}\left[ B \right]_{s,r}\left[ C \right]_{r,j}.
\]
Foglaljuk össze az eddigieket négyzetes mátrixok mellett:
\begin{proposition}
	Legyen $n\in\mathbb{N}$ természetes szám, és tekintsük az $n\times n$ méretű mátrixok
	halmazát, ellátva ezt a halmazt a mátrix összeadással és a mátrixszorzással.
	Az $\left( \mathbb{F}^{n\times n},+,\cdot \right)$ algebrai struktúra egy egységelemes gyűrű.
\end{proposition}
}
\frame{
Ez a gyűrű, az $n>1$ esetben biztosan nem kommutatív.
Például $n=2$ mellett
\[
	\begin{pmatrix}
		0 & 1 \\
		0 & 0
	\end{pmatrix}
	\cdot
	\begin{pmatrix}
		0 & 0 \\
		1 & 0
	\end{pmatrix}
	=
	\begin{pmatrix}
		1 & 0 \\
		0 & 0
	\end{pmatrix},
	\text{ amíg }
	\begin{pmatrix}
		0 & 0 \\
		1 & 0
	\end{pmatrix}
	\cdot
	\begin{pmatrix}
		0 & 1 \\
		0 & 0
	\end{pmatrix}
	=
	\begin{pmatrix}
		0 & 0 \\
		0 & 1
	\end{pmatrix}.
\]
Az sem igaz, hogy ez a gyűrű nullosztómentes lenne, hiszen például $n=2$ mellett az
\[
	A
	=
	\begin{pmatrix}
		0 & 1 \\
		0 & 0
	\end{pmatrix}
\]
mátrix nyilván nem a zérus mátrix (az összeadásra nézve neutrális elem),
de $A\cdot A=0$.

Ebből persze már az is következik, hogy a fenti $A$ mátrixnak nincs a szorzásra nézve inverze.
}
\frame{
A kommutativitás hiánya, az eddigiektől eltérő számolási gyakorlatot eredményez.
A számolás közben a mátrixok sorrendjén nem változtathatunk.
Persze előfordul, hogy két mátrix szorzata nem függ a sorrendtől.
Ilyenkor a két mátrixot egymással \emph{felcserélhetőnek}, vagy \emph{kommutálónak} mondjuk.\index{kommutáló mátrixok}

Például, az identitás mátrixszal minden más mátrix kommutál.
A mátrix azon elemeit, ahol a sor- és az oszlopindexek azonosak
\emph{fődiagonálisbeli}\index{fődiagonális}
elemeknek mondjuk.

Egy mátrixot
\emph{diagonális alakúnak}\index{diagonális mátrix}
mondunk,
ha minden nem zérus eleme a fődiagonálisában van.
Az is világos, hogy a diagonális mátrixok egymással kommutálnak.
}
\frame{
Nagyon is triviális mégis érdemes észrevenni, hogy a fentiek $n=1$ esetben nem jelentenek problémát.
Ilyenkor az $1\times 1$-es mátrixok tere voltaképpen azonos az $\mathbb{F}$-testtel, hiszen csak az a különbség,
hogy egy testbeli $a$ elemet $[a]$ módon írjuk. A szorzás és az összeadás definíciója ugyanazt adja,
ha mint a testbeli elemre, vagy az ebből képzett $1\times 1$-es mátrixra gondolunk.
}
\frame{
Az $n\times n$-es négyzetes mátrixok másik érdemleges részstruktúrája az
\[
	\mathcal{F}=
	\left\{ c\cdot I:c\in\mathbb{F} \right\}.
\]
Itt $I$ az $n\times n$ méretű identitás mátrix, tehát $\mathcal{F}$ elemei azon diagonális alakú mátrixok,
ahol minden elem a diagonálisban azonos.
Világos, hogy két ilyen mátrix összege és szorzata is ilyen marad:
\[
	aI+bI=\left( a+b \right)I\text{ és } aI\cdot bI=\left( ab \right)I.
\]
Ez azt jelenti, hogy az $\left( \mathcal{F},+,\cdot \right)$ struktúra egy egységelemes gyűrű.
Világos, hogy itt bármely két elem kommutál, ergo egy kommutatív egységelemes gyűrűvel állunk szemben és
az is teljesen nyilván való, hogy minden nem zérus elemnek van a szorzásra nézve inverze.
Azt kaptuk tehát, hogy a fenti $\mathcal{F}$ minden $n$ mellett egy test.
}
%\frame{\frametitle{Az inverz mátrix kétoldalisága}
%    Emlékezzünk arra, hogy ha egy mátrixon Gauss\,--\,Jordan-eliminációt hajtunk végre,
%    akkor az utolsó táblázatban vagy minden sor zérus, vagy minden sor tartalmaz pontosan egy eliminált változót.
%    \begin{proposition}
%        Ha két négyzetes mátrix szorzata az identitás mátrix,
%        akkor e két mátrix egymással kommutál.
%    \end{proposition}
%}
\section{A komplex számok mint mátrixok}
\frame{
\begin{definition}[Izomorf testek]\index{izomorfizmus}
	Legyenek $\mathbb{F}$ és $\mathbb{G}$ testek.
	Azt mondjuk, hogy a két test \emph{izomorf} egymással,
	ha létezik köztük \emph{művelettartó bijekció}, azaz létezik
	\[
		\varphi:\mathbb{F}\to\mathbb{G}
	\]
	bijekció, amely tartja a műveleteket is, azaz
	\[
		\varphi\left( a+b \right)=\varphi\left( a \right)+\varphi\left( b \right)
		\text{ és }
		\varphi\left( a\cdot b \right)=\varphi\left( a \right)\cdot \varphi\left( b \right).
	\]
	A művelettartó bijekciót \emph{izomorfizmusnak} nevezzük.
\end{definition}
}
\frame{
Az izomorf testek közt nem teszünk különbséget. Úgy tekintjük őket, hogy csak jelölésükben különböznek.
Például, ha az $\mathbb{R}$ valós számokra gondolunk,
akkor a $2\times 2$-es diagonális alakú valós mátrixok közül azok, ahol a diagonális mindkét eleme azonos,
a valós testtel izomorf testet alkot.
\[
	\mathcal{R}=\left\{
	\begin{pmatrix}
		a & 0 \\
		0 & a
	\end{pmatrix}
	:a\in\mathbb{R}
	\right\}
	\text{ és }
	\varphi:\mathbb{R}\to\mathcal{R},
	\text{ ahol }
	\varphi\left( a \right)
	=
	\begin{pmatrix}
		a & 0 \\
		0 & a
	\end{pmatrix}.
\]

A $2\times 2$-es valós mátrixok egységelemes gyűrűjében tehát $\mathcal{R}$ egy olyan részgyűrű, ami még test is,
és izomorf az $\mathbb{R}$ valós számtesttel.
\\
Voltaképpen a valós számtestet beágyaztuk a $2\times 2$-es mátrixok közé,
azaz egy $a$ valós számot az
\begin{math}
	\begin{pmatrix}
		a & 0 \\
		0 & a
	\end{pmatrix}
\end{math}
mátrixszal reprezentálunk (írunk le).

A $2\times 2$ méretű valós mátrixok még sok-sok más testet is tartalmaznak.
Ezek közül számunkra a legfontosabb a következő részhalmaz.
}
\frame{
\begin{defprop}[komplex számtest]
	Jelölje két tetszőleges $a,b\in\mathbb{R}$ valós szám mellett $M_{a,b}$ az
	$\left( a,b \right)$ valós számpárhoz tartozó
	\(
	M_{a,b}
	=
	\begin{pmatrix}
		a & -b \\
		b & a
	\end{pmatrix}
	\)
	mátrixot.
	Tekintsük az ilyen típusú mátrixok $\mathcal{C}$-vel jelölt halmazát:
	\[
		\mathcal{C}
		=
		\left\{
		M_{a,b}
		:a,b\in\mathbb{R}
		\right\}
	\]
	E részhalmaz
	\begin{enumerate}
		\item
		      zárt a mátrix összeadásra és a mátrix szorzásra,
		      így a $2\times 2$-es valós mátrixok egy speciális egységelemes részgyűrűje.
		\item
		      E részgyűrűben a mátrix szorzás kommutatív művelet,
		      és
		\item
		      és e részgyűrűben minden nem zérus mátrixnak van inverze is a mátrixszorzás műveletre nézve.
	\end{enumerate}
	Eszerint a $\left( \mathcal{C},+,\cdot \right)$ algebrai struktúra egy test.
	Ezt a testet nevezzük a \emph{komplex számtestnek}\index{komplex számok},
	vagy a \emph{komplex számtest mátrix reprezentációjának}.\index{komplex számok mátrix reprezentációja}
\end{defprop}
}
%\begin{proof}
%	Az $a,b,c,d\in\mathbb{R}$ valós számok mellett
%	\[
%		M_{a,b}+M_{c,d}=
%		\begin{pmatrix}
%			a & -b \\
%			b & a
%		\end{pmatrix}
%		+
%		\begin{pmatrix}
%			c & -d \\
%			d & c
%		\end{pmatrix}
%		=
%		\begin{pmatrix}
%			a+c & -b-d \\
%			b+d & a+c
%		\end{pmatrix}
%		=
%		M_{a+c,b+d}
%	\]
%	és hasonlóan a mátrix szorzás definíciója szerint
%	\[
%		M_{a,b}\cdot M_{c,d}=
%		\begin{pmatrix}
%			a & -b \\
%			b & a
%		\end{pmatrix}
%		\cdot
%		\begin{pmatrix}
%			c & -d \\
%			d & c
%		\end{pmatrix}
%		=
%		\begin{pmatrix}
%			ac-bd & -ad-bc \\
%			bc+ad & -bd+ac
%		\end{pmatrix}
%		=
%		M_{ac-bd,ad+bc}.
%	\]
%	Mivel $M_{1,0}$ a $2\times 2$-es identitás mátrix, ezért $\mathcal{C}$ a mátrix összeadásra és a mátrixszorzásra nézve
%	egységelemes gyűrűt alkot.
%
%	A szorzás kommutativitása is látszik a fenti számolásból, hiszen
%	\[
%		M_{c,d}\cdot M_{a,b}=M_{ca-db,cb+da}=M_{ac-bd,ad+bc}=M_{a,b}\cdot M_{c,d}
%	\]
%	a valós számok összeadásának és szorzásának
%	kommutativitása miatt.
%
%	Legyen most $M_{a,b}$ egy nem zérus mátrix, így $a^2+b^2\neq 0$.
%	Világos, hogy
%	\[
%		M_{a,b}\cdot M_{a,-b}=M_{a^2+b^2,0}=\left( a^2+b^2 \right)M_{1,0}=\left( a^2+b^2 \right)I.
%	\]
%	Ebből már látszik is, hogy $M_{a,b}\cdot M_{\frac{a}{a^2+b^2},\frac{-b}{a^2+b^2}}=I$.
%	Ez a már igazolt kommutativitással éppen azt jelenti, hogy minden nem zérus elemnek van multiplikatív inverze,
%	ergo $\mathcal{C}$ valóban test.
%\end{proof}
\frame{
Ebben a $\mathcal{C}$ testben az
\(
M_{0,1}=
\begin{pmatrix}
	0 & -1 \\
	1 & 0
\end{pmatrix}
\)
elem olyan, hogy a négyzete a szorzásra nézve reprodukáló elemnek az összeadásra nézve képzett inverze:
\[
	M_{0,1}\cdot M_{0,1}
	=
	\begin{pmatrix}
		0 & -1 \\
		1 & 0
	\end{pmatrix}
	\cdot
	\begin{pmatrix}
		0 & -1 \\
		1 & 0
	\end{pmatrix}
	=
	\begin{pmatrix}
		-1 & 0  \\
		0  & -1
	\end{pmatrix}
	=
	M_{-1,0}
	=
	-I.
\]
Ez a tulajdonság azért figyelemre méltó,
mert ha a szokásoknak megfelelően a test multiplikatív neutrális elemét az $1$ szimbólummal jelöljük,
akkor olyan elemet találtunk a komplex számtestben,
amelynek négyzete éppen $-1$.
Tudjuk, hogy a valós számtest esetében ez nem lenne lehetséges.
}
\frame{
Tekintsük most valamely $a,b\in\mathbb{R}$ mellett az
\[
	M_{a,b}=
	\begin{pmatrix}
		a & -b \\
		b & a
	\end{pmatrix}
	=
	\begin{pmatrix}
		a & 0 \\
		0 & a
	\end{pmatrix}
	+
	\begin{pmatrix}
		0 & -1 \\
		1 & 0
	\end{pmatrix}
	\cdot
	\begin{pmatrix}
		b & 0 \\
		0 & b
	\end{pmatrix}
\]
felbontást.
Ha bevezetjük az
\(
i=M_{0,1}=
\begin{pmatrix}
	0 & -1 \\
	1 & 0
\end{pmatrix}
\)
jelölést\index{i komplex szám@$i$ komplex szám}, akkor minden komplex szám
\[
	M_{a,b}=
	\begin{pmatrix}
		a & 0 \\
		0 & a
	\end{pmatrix}
	+
	i
	\cdot
	\begin{pmatrix}
		b & 0 \\
		0 & b
	\end{pmatrix}
\]
alakban írható.
}
\frame{
Emlékezzünk arra,
hogy a valós számtest is részhalmaza a komplex számtestnek abban az értelemben,
ha minden $a$ valós számot az
\begin{math}
	\begin{pmatrix}
		a & 0 \\
		0 & a
	\end{pmatrix}
\end{math}
mátrixszal reprezentálunk. ($\mathcal{R}\subseteq\mathcal{C}$).
Ha tehát megegyezünk abban, hogy az $a$ valós számra nézve mi az
\(
\begin{pmatrix}
	a & 0 \\
	0 & a
\end{pmatrix}
\)
mátrixra gondolunk,
akkor azt kapjuk, hogy minden komplex szám
\[
	M_{a,b}=a+ib
\]
alakú, ahol $i$ egy olyan komplex szám, amelyre $i^2=-1$, $a,b\in\mathbb{R}$.
Ezt nevezzük a komplex szám \emph{normálalakjának}.\index{komplex szám normálalakja}
}
\frame{
Ne felejtsük a műveleteket:
Láttuk, hogy $M_{a,b}+M_{ c,d}=M_{a+c,b+d}$.
Ez a normálalak reprezentáció mellett azt jelenti, hogy az összeadás definíciója csak
\metroset{block=fill}\begin{block}{}
\[
	(a+ib)+(c+id)=(a+c)+i(b+d)\tag{\dag}
\]
\end{block}\metroset{block=transparent}
lehet.
Hasonlóan emlékszünk, hogy
\begin{math}
	M_{a,b}\cdot M_{c,d}=M_{ac-bd,ad+bc}
\end{math},
ami normálalak reprezentáció mellett az
\metroset{block=fill}\begin{block}{}
\[
	(a+ib)\cdot (c+id)=(ac-bd)+i(ad+bc)\tag{\ddag}
\]
\end{block}\metroset{block=transparent}
definíciót eredményezi.
}
\frame{
A következőket gondoltuk meg:
\begin{defprop}[komplex számtest a normálakkal]
    Jelölje
    \metroset{block=fill}
    \begin{alertblock}{}
	\[
		\mathbb{C}
		=
		\left\{
		a+ib:a,b\in\mathbb{R}
		\right\}
	\]
\end{alertblock}\metroset{block=transparent}
	a \emph{komplex számok normálalakját}.
	Az összeadás műveletet definiálja (\dag), és a szorzás műveletet definiálja (\ddag).
	Az így kapott algebrai struktúra test, amely izomorf a komplex számtest mátrix reprezentációjával.
	Az izomorfizmust a
	\[
		\varphi:\mathcal{C}\to\mathbb{C},\quad
		\varphi
		\left(
		\begin{matrix}
				a & -b \\
				b & a
			\end{matrix}
		\right)
		=
		a+ib
	\]
	művelettartó bijekció hozza létre.
\end{defprop}
}
\frame{
Ha már megértettük, hogy a komplex számok normálalak reprezentációja testet alkot,
akkor a (\dag) és (\ddag) definíciók megjegyzése nagyon könnyű.
Más nem is lehet:
Az összeadáshoz (\dag) csak el kell végezni a műveletet majd kiemelni $i$-t,
a szorzás definíciójához (\ddag)
az $i$ kiemelése után jutunk.
}

\section{A komplex számok abszolútértéke}
\frame{
A valós számokra jól ismert abszolútérték függvényt terjesztjük ki komplex számtest elemeire.
\begin{definition}[valós rész, képzetes rész]\index{komplex szám valós része}\index{komplex szám képzetes része}
	Legyen $z\in\mathbb{C}$, $z=a+ib$.
	Ekkor $a\in\mathbb{R}$ a $z$ komplex szám \emph{valós része},
	és $b\in\mathbb{R}$ a $z$ komplex szám \emph{képzetes része}.
	$\Re z$ jelöli a valós részt, és $\Im z$ a képzetes részt.
\end{definition}
\begin{definition}[konjugált]\index{komplex szám konjugáltja}
	Legyen $z\in\mathbb{C}$, $z=a+ib$.
	Ekkor $z$ \emph{konjugáltja} $\bar{z}=a-ib$.
\end{definition}
}
\frame{
A konjugált definíciója alapján, egyszerű számolgatással kapjuk az alábbi azonosságokat.
\begin{proposition}
	Minden $z\in\mathbb{C}$ komplex szám mellett fennállnak a konjugálás következő szabályai:
	\[
		\overline{\left( \bar{z} \right)}=z,\quad
		z+\bar{z}=2\Re z,\quad
		z-\bar{z}=2i\Im z,\quad
		z\in\mathbb{R}\iff z=\bar{z},\quad
		z\bar{z}=(\Re z)^2+(\Im z)^2\geq 0.
	\]
	Az összeadás és szorzás művelet is felcserélhető a konjugálással,
	azaz bármely két $z,w\in\mathbb{C}$ komplex szám esetén igaz, hogy
	\[
		\overline{z+w}=\bar{z}+\bar{w},\quad
		\overline{z\cdot w}=\bar{z}\cdot\bar{w},\quad
		\overline{z-w}=\bar{z}-\bar{w},\quad
		\overline{\left( \frac{z}{w} \right)}=\frac{\bar{z}}{\bar{w}}\quad\text{ feltéve, hogy }w\neq 0.\qedhere
	\]
\end{proposition}
}
\frame{
Most használjuk először a valós számok rendezését.
Az $\mathbb{R}$ testen a $\geq $ reflexív, antiszimmetrikus, tranzitív relációt az algebrai műveletekkel a következő két axióma kapcsolja össze:
Minden $a,b,c\in\mathbb{R}$ mellett
\[
	a\geq b\text{ esetén }a+c\geq b+c,\qquad a,b\geq 0\text{ esetén }ab\geq 0.
\]
Az $a^2-b^2=\left( a+b \right)\left( a-b \right)$ azonosság szerint,
az $a,b$ nem negatív valós számok mellett $a\geq b$ akkor és csak akkor, ha $a^2\geq b^2$.
Speciálisan $a^2=b^2$ akkor és csak akkor, ha $a=b$.
}
\frame{
\begin{definition}[komplex szám abszolútértéke]\index{abszolútérték}
	Legyen $z\in\mathbb{C}$ egy komplex szám.
	Láttuk, hogy
	$z\bar{z}\in\mathbb{R}$ és
	$z\bar{z}\geq 0$.
	E szám négyzetgyökét nevezzük a $z$ komplex szám \emph{abszolútértékének}.
	Jelölés: $|z|=\sqrt{z\bar{z}}$.
\end{definition}
Ha speciálisan $\Im z=0$,
tehát ha $z$ egy valós szám,
akkor e definíció szerint
\[
	|z|=
	\sqrt{z^2}=
	\begin{cases}
		z  & \text{, ha }z\geq 0     \\
		-z & \text{, ha }z<0\text{,}
	\end{cases}
\]
ami egybeesik a valós számok abszolútértékének definíciójával.
A fentiből azonnal látszik, hogy minden $a\in\mathbb{R}$ valós szám mellett $a\leq|a|$.
}
\frame{
Mivel $z\bar{z}=\bar{z}z$, ezért $|z|=|\bar{z}|$, azaz komplex számnak és konjugáltjának azonos az abszolútértéke.
Szokásos technika, hogy a komplex szám abszolútértékére vonatkozó állításokat, az abszolútértékek négyzetére fogalmazzuk át.
Ez megtehető, hiszen nem negatív valós számok egyenlősége és azok négyzetének egyenlősége azonos fogalmak.
Például $|z|=0$, akkor és csak akkor, ha $z\bar{z}=0$, ami akkor és csak akkor teljesül ha $z=0$.
Hasonlóan, ha $z,w$ komplex számok, akkor $|zw|^2=zw(\overline{zw})=zw\bar{z}\bar{w}=z\bar{z}w\bar{w}=|z|^2|w|^2=\left( |z||w| \right)^2$,
ergo $|zw|=|z||w|$.

Világos, hogy $|\Re z|^2\leq(\Re z)^2+(\Im z)^2=z\bar{z}=|z|^2$, emiatt
\(
|\Re z|\leq |z|.
\)
Ezt alkalmazva tetszőleges $z,w$ komplex számok mellett
\[
	|
	\Re (z\bar{w})
	|
	\leq
	|z\bar{w}|
	=
	|z||w|,
\]
amit \emph{Cauchy\,--\,Schwartz-egyenlőtlenségnek}\index{Schwartz-egyenlőtlenség} mondunk.
}
\frame{\frametitle{Az abszolútérték legfontosabb tulajdonságai}
\metroset{block=fill}
\begin{proposition}
	Legyen $z,w\in\mathbb{C}$ komplex szám.
	Ekkor
	\begin{enumerate}
		\item $|z|=0$ akkor és csak akkor, ha $z=0$,
		\item $|zw|=|z||w|$,
		\item $|z+w|\leq |z|+|w|$.\qedhere
	\end{enumerate}
\end{proposition}
\metroset{block=transparent}
Az utolsó egyenlőtlenséget \emph{háromszög-egyenlőtlenségnek}\index{háromszög-egyenlőtlenség} nevezik.
Egy kevésbé népszerű, de ekvivalens alakja
\[
	\left|\left( |z|-|w| \right)\right|\leq|z-w|.
\]
}
%\begin{proof}
%	Csak a háromszög-egyenlőtlenséget nem gondoltuk meg eddig:
%	\begin{multline*}
%		|z+w|^2=
%		\left( z+w \right)\left( \overline{z+w} \right)
%		=
%		\left( z+w \right)\left( \bar{z}+\bar{w} \right)
%		=
%		z\bar{z}+w\bar{w}+z\bar{w}+\overline{z\bar{w}}
%		=
%		\\
%		|z|^2+|w|^2+2\Re (z\bar{w})
%		\leq
%		|z|^2+|w|^2+2|\Re (z\bar{w})|
%		\leq
%		|z|^2+|w|^2+2|z||w|
%		=
%		\left( |z|+|w| \right)^2.
%	\end{multline*}
%	Az első egyenlőtlenség, hogy egy valós szám nem nagyobb az abszolútértékénél,
%	a második egyenlőtlenség pedig a Cauchy\,--\,Schwartz-egyenlőtlenség.
%
%	A kevésbé populáris alakhoz használjuk a már igazolt egyenlőtlenséget:
%	\begin{math}
%		|z|
%		=
%		|(z-w)+w|
%		\leq
%		|z-w|+|w|,
%	\end{math}
%	majd ugyanezt a $z$ és a $w$ számok felcserélése után mégegyszer felírva
%	\[
%		|z|-|w|\leq |z-w|
%		\quad\text{ és }\quad
%		|w|-|z|\leq |w-z|.
%	\]
%	Na most, azt ugyan nem tudjuk, hogy a $||z|-|w||$ szám a baloldali számok közül melyikkel azonos,
%	de az egyikkel biztosan egybeesik.
%	No de, mindegy is melyikkel, hiszen a jobboldal mindkét esetben a kívánt $|z-w|$.
%\end{proof}

\section{A komplex számok trigonometrikus alakja}
\frame{
Láttuk, hogy $z,w\in\mathbb{C}$ komplex szám mellett
\[
	\Re\left( z+w \right)=\Re z+\Re w
	\qquad
	\text{ és }
	\qquad
	\Im\left( z+w \right)=\Im z+\Im w.
\]
Ez azt jelenti, hogy ha az $a+ib$ komplex számot azonosítjuk az $\mathbb{R}^2$ sík
$\left( a,b \right)$ pontjával,
akkor egyszerűen koordinátánként kell összeadni a komplex számokat, mintha a $z$ komplex szám
az origóból az $\left( a,b \right)$ pontra mutató vektor lenne.

A kérdés, hogy ha így képzeljük a komplex számokat,
akkor \alert{\bf komplex számok szorzása mit jelent a komplex számoknak megfeleltetett vektorok körében?}
}
\frame{
\begin{defprop}[komplex szám trigonometrikus alakja]\index{komplex szám trigonometrikus alakja}
	Legyen $z\in \mathbb{C}$ egy nem zérus komplex szám.
	Ekkor létezik $\varphi\in\mathbb{R}$ valós szám, hogy
	\[
		z=
		|z|\left( \cos\varphi+i\sin\varphi \right)\tag{\dag}
	\]
	Egy ilyen $\varphi$ számot nevezzük a $z$ komplex szám egy \emph{argumentumának}, és néha $\arg z$ módon jelöljük.
	\index{komplex szám argumentuma}
	A (\dag) alak a komplex szám \emph{trigonometrikus alakja}.

	Két nem zérus komplex szám egyenlősége azt jelenti, hogy az abszolútértékük azonos,
	és az argumentumaik különbsége $2\pi$ többszöröse.
\end{defprop}
}
%\begin{proof}
%	Világos, hogy ha $z=a+ib\neq 0$, akkor $a^2+b^2>0$, így
%	\[
%		z=a+ib
%		=
%		\sqrt{a^2+b^2}\left( \frac{a}{\sqrt{a^2+b^2}}+i\frac{b}{\sqrt{a^2+b^2}} \right).
%	\]
%	Ha $x$ jelöli a fenti zárójelben a valós részt és $y$ a képzetes részt,
%	akkor $x^2+y^2=1$.
%	Így az $\left( x,y \right)$ pár a sík egységkörének egy pontja.
%	A trigonometrikus függvények középiskolai definíciója szerint,
%	ha $\varphi$ jelöli az $\left( 1,0 \right)$ pontot az $\left( x,y \right)$
%	ponttal összekötő ív hosszát, -- az óramutató járásával ellentétes irányban mérve a körcikk peremén --
%	akkor $x=\cos\varphi$ és $y=\sin\varphi$.
%\end{proof}
\frame{
\begin{proposition}
	Legyenek a $z,w\in \mathbb{C}$ nem nulla komplex számok a trigonometrikus alakjukban felírva,
	azaz
	\[
		z=|z|\left( \cos\varphi+i\sin\varphi \right),\qquad
		w=|w|\left( \cos\psi+i\sin\psi \right).
	\]
	Ekkor a két szám szorzatának abszolútértéke az abszolútértékek szorzata
	és a szorzat egy argumentumát is megkapjuk mint az argumentumok összege.
	Magyarul:
	\[
		zw=
		|z||w|\left( \cos\left( \varphi+\psi \right)+i\sin\left( \varphi+\psi \right) \right).
	\]
	Speciálisan minden $n\in\mathbb{Z}$ egész számra\index{Moivre-formula}
	\[
		z^n=|z|^n\left( \cos n\varphi+i\sin n\varphi \right).\qedhere
	\]
\end{proposition}
}
\frame{
A szakasz bevetőjében feltett kérdésre tehát a válasz,
hogy a $z$ komplex számmal való szorzást a sík olyan geometriai transzformációjának képzelhetjük,
amely egy $\arg z$ szöggel való forgatásból és egy $|z|$-szeres origó középpontú nyújtásból áll.
}
\frame{
\begin{definition}[egységgyök]\index{komplex $n$-edik egységgyökök}
	Legyen az $n\in\mathbb{N}$ természetes szám rögzítve.
	Tetszőleges $k=0,1,\dots,n-1$ mellett jelölje
	\[
		\epsilon_k^{(n)}=\cos k\frac{2\pi}{n}+i\sin k\frac{2\pi}{n}.
	\]
	az úgynevezett \emph{$n$-edik komplex egységgyököket}.
\end{definition}
Az $n$-edik komplex egységgyökök a komplex számsík pontosan $n$ különböző pontját alkotják,
és az $n$-edik hatványuk $1$, azaz
$(\epsilon_k^{(n)})^n=1$.
Emiatt minden $w\neq 0$ komplex szám mellett pontosan $n$ komplex gyöke van a $t^n-w$ polinomnak.
Ha ugyanis $w$ trigonometrikus alakja $w=|w|\left( \cos\psi+i\sin\psi \right)$,
akkor legyen például $z_0=\sqrt[n]{|w|}\left( \cos(\psi/n)+i\sin(\psi/n) \right)$,
és így $(z_0\epsilon_k^{(n)})^n=w$ minden $k=0,\dots,n-1$ szám mellett.
}
\end{document}
